\documentclass[10pt]{article}

\usepackage[utf8]{inputenc}

\usepackage[T1]{fontenc}

\usepackage{amsmath}

\usepackage{amsfonts}

\usepackage{amssymb}

\usepackage[version=4]{mhchem}

\usepackage{stmaryrd}

\usepackage{graphicx}

\usepackage[export]{adjustbox}

\graphicspath{ {./images/} }



\begin{document}

при условии, что этот предел существует для всех $\varepsilon$, $0<\varepsilon<\pi / 2$, и, следовательно, не зависит от $\varepsilon$. Этот термин иногда употребляется в обобщенном смысле для функций $f(z)$, заданных в произвольных (включая многомерные) областях $D$, причем в качестве $S(\zeta, \varepsilon)$ берется пересечение с $D$ угловой (или конической) области с вершиной $\zeta \in \partial D$, осью к-рой является нормаль к границе $\partial D$ в точке $\zeta$, а угол раствора равен $\pi / 2-\varepsilon$, $0<\varepsilon<\pi / 2$.



Лит. [1] м а р к у ш е в и ч А. И., Теория аналитических функций, 2 изд., т. 1-2, М , 1967-68, [2] ПГ р и в а л о в И. И., Граничные свойства аналитических функций, $\underset{E}{2}$ изд., М С-Л,

1950.



УГОЈI - геометрическая фигура, состоящая из двух различных лучей, выходящих из одной точки. Лучи наз. с т о р о н а м и У., а их общее начало - в е pші и но й У. Пусть $[B A),[B C)$ - стороны угла, $B$ - его вершина, $\alpha$ - плоскость, определяемая сторонами У. Фигура $\Gamma=[B A) \cup[B C)$ делит плоскость $\alpha$ на две фипуры $\alpha_{i}^{\prime}=\alpha \backslash \Gamma, i=1,2: \alpha_{1}^{\prime} \bigcup \alpha_{2}^{\prime}=\alpha \backslash \Gamma$. Фигура $\alpha_{i}=\alpha_{i}^{\prime} \cup \Gamma, i=1,2$, также наз. У. или п ло о с к и м у г ло м, $\alpha_{i}^{\prime}$ наз. внутренней областью плоскоro $\mathbf{Y}, \alpha_{i}$.



Два угла наз. р а в ны м и (или к о нг р у ә т тн ы м и), если они могут быть совмещены так, что совпадут их соответствующие стороны и верпины. От любого луча на плоскости в данную сторону от него можно отложить единственный $\mathrm{Y}$., равный данному $\mathrm{y}$. Сравнение У. осуществляется двумя способами. Если



\begin{center}

\includegraphics[max width=\textwidth]{2023_07_09_b2b43e0bf64f9e28e51fg-1(1)}

\end{center}



Рис. 1. У. рассматривается как пара лучей с общим началом, то для выяснения вопроса, какой из двух У. больше, необходимо совместить в одной плоскости вершины У. и одну пару их сторон (см. рис. 1). Если вторая сторона одного У. окажется расположенной внутри другого У., то говорят, что первый У. меньше, чем второй. В'торой способ сравнения $У$. основан на сопоставлении каждому У. нек-рого числа. Равным У. будет соответствовать одинаковое число градусов или радиан (cM. ниже), большему У.- большее число, меньшему меньшее.



Два У. наз. с м е ж н ы м и, если у них общая вершина и одна сторона, а две другие стороны образугт



\begin{center}

\includegraphics[max width=\textwidth]{2023_07_09_b2b43e0bf64f9e28e51fg-1}

\end{center}



Рис. 2. прямую (см. рис. 2). Вообще, У., имеющие общую вершину и одну общую сторону, наз. ш р и ле жа и и и. У. наз. в е р т и а ль н ы и, если стороны одного являются продолжениями за вершину сторон другого У. Вертикальные У. равны между собой. У., у к-рого стороны образуют прямую, наз. р а з в е р н у т ы м. Половина развернутого У. наз. п р я м ы м У. Прямой У. можно әквивалентно определить иначе: У., равный своему смежному, наз. п р я м ы м. Внутренняя область плоского У., не превосходящего развернутого, является выпуклой областью на плоскости.



За единицу измерения $У$. принимается 90 -я доля шрямого У., наз. градусом.



Используется и т.н. к р уговая, или р а д иа н н а я, мера У. Числовое значение радианной меры У. равно длине дуги, высекаемой сторонами У. из единичной окружности. Один радиан приписывается $\mathrm{У}$., соответствующему дуге, длина к-рой равна ее радиусу. Развернутый У. равен $\pi$ радиан.



При пересечении двух прямых, лежащих в одной плоскости, третьей прямой образуются У. (см. рис. 3 ): 1 и 5,2 и 6,4 и 8,3 и 7 - наз. с о о т в е т с т в н-



ным и; 2 и 5,3 и $8-$ вну т ренними однос тор онними; 1 и 6,4 и 7 - вне шними о д носторон ними; 3 и 5,2 и $8-$ внутренними накре с т ле жащ ии; 1 и 7 , 4 и 6 - в е шн и ми нак р е с т ле жа ш и и.



В практич. задачах целесообразно рассматривать У. как меру ліоворота фиксированного луча вокруг его начала до заданного положения. В зависимости от направления Ііоворота У. в этом случае можно рассматривать как по-



\begin{center}

\includegraphics[max width=\textwidth]{2023_07_09_b2b43e0bf64f9e28e51fg-1(2)}

\end{center}



Рис. 3 .



ные. Тем самым У . в этом смысле может иметь своей величиной любое действительное число. У. как мера поворота луча рассматривается в теории тригонометрич. функций: для любых значений аргумента (У.) можно определить значения тригонометрич. функций. Понятие У. в геометрич. системе, в основу к-рой положена точечно-векторная аксиоматика, в корне отличается от определений У. как фигуры - в этой аксиоматике под $\mathbf{y}$. понимают определенную метрич. величину, связанную с двумя векторами с помощьто операции скаля рного умножения векторов. Именно, каждая пара векторов $\boldsymbol{a}$ и $\boldsymbol{b}$ определяет нек-рый угол $\varphi-$ число, связанное с векторами формулой



\[

\cos \varphi=\frac{(a, b)}{|a||b|}

\]



где $(\boldsymbol{a}, \boldsymbol{b})$ - скалярное произведение векторов.



Понятие У. как плоской фигуры и как нек-рой числовой величины применяется в различных геометрич. задачах, в к-рых У. определяется специальным образом. Так, под У. между пересекающимися кривыми, имеющими определенные касательные в точке пересечения, понимают У., образованный ртими касательными.



За угол между прямой и плоскостью принимается У., образованный прямой и ее прямоугольной проекцией на плоскость; он измеряется в пределах от $0^{\circ}$ ло $90^{\circ}$. Под У. между скрещиваюпимися прямыми понимают У. между направлениями этих прямых, т. е. между прямыми, параллельными, скрещивающимися и проведенными через одну точку.



Телесным уәлом наз. часть пространства, ограниченная нек-рой конической поверхностью; частным случаем телесного У. является многогранышй угол.



В многомерной геометрии определяется $\mathrm{Y}$. между многомерными плоскостями, между прямыми п плоскостями и т. д. Углы между прямыми, плоскостями и прямыми, многомерными плоскостями определяются и в неевклидовых пространствах. J $A$. Cuдоров.



УДАРНЫХ ВОЛН МАТЕМАТИЧЕСКАЯ ТЕОРИЯ математическое описание свойств, движения и взаимодействия с окружающей средой поверхностей разрыва параметров среды (ударных волн). В более широком и абстрактном смысле У. в. м. т. огисывает свойства поверхностей разрыва решениї квазилинейных аиперболических уравнений и систем. У. в. М. т. вовникла в связи с задачами двинения газов и сжимаемых жидкостей во 2-й пол. 19 в.; ее основы были заложены в работах С. Ирншоу, Б. Римана, У. Ранкина, П. Гюгоньо (см., напр., [1] - [4]).



При идеализации реальных газов и жидкостей рассматривают среды, лишенные диссипативных свойств, в к-рых отсутствуют вязкость и теплопроводность. В процессе движения в таких идеальных средах могут возникать разрывы в распределениях всех параметров течения (плотности, давления, температуры, скорости и др.). Множества точек разрыва параметров течения





\end{document}